\documentclass[11pt]{letter}
\usepackage[margin=1in]{geometry}
\usepackage{times}
\usepackage[hidelinks]{hyperref}
\signature{Jaiveer Bassi\\Independent Researcher\\\href{mailto:jaiveerbassi@yahoo.com}{jaiveerbassi@yahoo.com}}
\address{Jaiveer Bassi\\Independent Researcher, USA\\\href{mailto:jaiveerbassi@yahoo.com}{jaiveerbassi@yahoo.com}}
\date{\today}

\begin{document}
\begin{letter}{Editors\\\textit{Scientific Reports}}
\opening{Dear Editors,}

Please consider the Article ``Complete Patient-Level Leakage Among Traceable Images in a Widely Used Brain Tumor MRI Benchmark'' for publication in \textit{Scientific Reports}.

The manuscript audits a widely used four-class brain MRI benchmark whose aggregated release discarded patient identity. By matching 4,586 tumor images back to their figshare source, I recover subject identifiers and show that every traceable tumor test image belongs to a patient represented in training. The paper then re-evaluates the same pipeline under patient-disjoint cross-validation and includes a matched random-fold control over the identical 6,292-image population. The control attributes 2.73 of the 4.03 percentage-point performance gap to patient grouping; the class without patient identifiers changed negligibly despite a smaller relaxation from near-duplicate to exact-pixel grouping. The work also distinguishes patient leakage from exact-image duplication, which does not explain the near-ceiling result.

The study is appropriate for \textit{Scientific Reports} because it presents a technically validated correction to a public medical-imaging benchmark, quantifies the effect with paired and grouped analyses, and releases the patient map, immutable folds, predictions, and complete analysis code. The finding affects the interpretation and reproducibility of a broad body of transfer-learning work while remaining carefully limited to this dataset and pipeline.

This manuscript is not under consideration elsewhere. An earlier preprint by the same author used the released split and predates the audit; it is identified in the submission and is superseded by this work. I have had no prior discussions about this manuscript with a \textit{Scientific Reports} Editorial Board Member. I have no competing interests and no reviewers to exclude. The work is a secondary analysis of public, de-identified data. AI assistance is disclosed in the manuscript, and I take full responsibility for the study and submitted materials.

I suggest the following reviewers, none of whom I have collaborated with, been supervised by, or corresponded with about this work:

\begin{itemize}
\itemsep2pt
\item Prof.\ Stefano Diciotti, Department of Electrical, Electronic and Information Engineering ``Guglielmo Marconi'', University of Bologna, Italy (\texttt{stefano.diciotti@unibo.it}) --- senior author of the closest precedent for slice-level leakage in brain MRI classification.
\item Prof.\ Luca Citi, School of Computer Science and Electronic Engineering, University of Essex, United Kingdom (\texttt{lciti@essex.ac.uk}) --- co-author of the same study, at a separate institution.
\item Dr.\ Lauren Oakden-Rayner, Australian Institute for Machine Learning, University of Adelaide, Australia (\texttt{lauren.oakden-rayner@adelaide.edu.au}) --- radiologist and machine-learning researcher whose work audits the labelling and construction of public medical-imaging datasets.
\item Prof.\ Jayashree Kalpathy-Cramer, Division of Artificial Medical Intelligence, Department of Ophthalmology, University of Colorado Anschutz School of Medicine, USA (\texttt{jayashree.kalpathy-cramer@cuanschutz.edu}) --- works on evaluation rigour and reproducibility in clinical imaging AI, including partitioning of radiology datasets.
\end{itemize}

The exact manuscript code and artifacts are public in GitHub release v1.2.2 and the corresponding Zenodo archive (DOI: 10.5281/zenodo.22755070).

Thank you for your consideration.

\closing{Sincerely,}
\end{letter}
\end{document}
